% Corrected proof blocks only; this is not the Paper II manuscript.
% Exact source IDs and existing labels are preserved.
% Checker: GPT-6 Astra (OpenAI, ChatGPT Pro mode), effort Pro, writing run,
% 17 November 2026 UTC.
% For A.1.1, S is a reduced Cartier hypersurface of a smooth complex
% manifold Y; N^*=I_S/I_S^2, V=Omega_Y^1|_S,
% T=tors(Omega_S^1), Q=Omega_S^1/T, and A is a line bundle on S.
% These definitions must be stated in prose before the block in the paper.

\begin{statement}{A.1.1}{Proved}{The conormal filtration and its exact section criterion}
The conormal map $d:N^*\to V$, defined by $[F]\mapsto dF|_S$ for a local defining equation $F$, is injective. With $K=\ker(V\to Q)$, there are canonical exact sequences
\begin{align}
0&\longrightarrow N^*\xrightarrow{d}K\longrightarrow T\longrightarrow0,
\label{A:filtration1}\\
0&\longrightarrow K\longrightarrow V\longrightarrow Q\longrightarrow0.
\label{A:filtration2}
\end{align}
Let $\partial_A:H^0(S,Q\otimes A)\to H^1(S,K\otimes A)$ and $\delta_A:H^0(S,T\otimes A)\to H^1(S,N^*\otimes A)$ be their connecting maps. Then
\[
 H^0(S,V\otimes A)=0
 \quad\Longleftrightarrow\quad
 H^0(S,K\otimes A)=0\ \hbox{and}\ \partial_A\hbox{ is injective}.
\tag{A.1}
\]
Under the explicit hypotheses $H^0(S,N^*\otimes A)=0$ and $H^0(S,Q\otimes A)=0$, $H^0(S,V\otimes A)$ is canonically $\ker\delta_A$.
\end{statement}
\begin{proof}
Locally $R=\C\{x_0,\ldots,x_n\}/(F)$ is reduced. Its total quotient ring is the product of the fraction fields of its irreducible components. At the generic point of every component, $S$ is smooth and $dF$ is nonzero. Consequently $a\,dF=0$ in $R^{n+1}$ implies that $a$ vanishes in each of those fields. The map from the reduced ring $R$ into their product is injective, so $a=0$. This proves injectivity in the usual right-exact conormal sequence. Taking the inverse image of $T\subset\Omega_S^1$ gives \eqref{A:filtration1} and \eqref{A:filtration2}. Tensoring by $A$ is exact because $A$ is locally free. The cohomology sequence of \eqref{A:filtration2} identifies the image of $H^0(VA)$ in $H^0(QA)$ with $\ker\partial_A$ and its kernel with $H^0(KA)$, proving (A.1). When $H^0(QA)=0$, it identifies $H^0(VA)$ with $H^0(KA)$. The first sequence then identifies this group with $\ker\delta_A$ when $H^0(N^*A)=0$.
\end{proof}

\begin{statement}{A.2.6}{Proved}{A transverse cusp whose critical-curve equation is only a symbolic square}
Put
\[
 \begin{aligned}
 I&=(y^2-xz,\ yz^2-x^4,\ z^3-x^3y)\subset\C\{x,y,z\},\\
 F&=4x^7-15x^3yz^2+10x^2y^3z-3xy^5+4z^5.
 \end{aligned}                                                   \tag{A.NL.1}
\]
The reduced critical locus of the reduced hypersurface $F=0$ is the
irreducible curve $D$ parametrized by $(x,y,z)=(t^5,t^6,t^7)$.
Its transverse type at every point with $t\ne0$ is the ordinary cusp
$A_2$. The curve is not a local complete intersection at the origin,
and
\[
             F\in I^{(2)}\setminus I^2.                         \tag{A.NL.2}
\]
Thus the ordinary-square membership used to define the conormal
quadratic jet in A.5.5 does not extend to an arbitrary reduced
critical curve, even with generic transverse type $A_2$.
This is a local example, without an asserted compact CDP realization.
\end{statement}
\begin{proof}
Write the three displayed generators of $I$ as $g_1,g_2,g_3$.
The quotient is a finite $\C\{x\}$-module generated by
$1,y,z,yz,z^2$: their products by $y$ and $z$ reduce to this list
using $y^2=xz$, $yz^2=x^4$, and $z^3=x^3y$.
For example these identities also give $y^5=x^6$ and $z^5=x^7$,
which make the analytic quotient finite over $\C\{x\}$.
The parametrization sends the five generators to
$1,t^6,t^7,t^{13},t^{14}$. Their exponents have distinct residues
modulo five, so they are linearly independent over $\C\{t^5\}$.
Consequently this actual ring homomorphism is injective and has
image $\C\{t^5,t^6,t^7\}$; it proves the stated normalization and
the primeness of $I$. The three generators are minimal. Indeed their
initial terms have degrees two, three, three; after the degree-two
coefficient vanishes, the cubic terms $yz^2$ and $z^3$ are independent
modulo the span of $xg_1,yg_1,zg_1$. A height-two ideal with three
minimal generators is not a complete intersection.

The following polynomial identities determine the critical locus:
\[
 \begin{aligned}
 F_y&=-15xg_1^2,\\
 F_x&=(17xyz-3y^3)g_1-28x^2g_2,\\
 F_z&=10x^2yg_1+20zg_3.
 \end{aligned}                                                   \tag{A.NL.3}
\]
At $x=0$, the equations $F_x=F_z=0$ give $y=z=0$.
At $x\ne0$, $F_y=0$ gives $g_1=0$, and then $F_x=0$
gives $g_2=0$. The value $z=0$ is impossible there, so $F_z=0$
also gives $g_3=0$. Conversely all three partials vanish on $D$.
Weighted Euler's identity, with weights $(5,6,7)$ and degree $35$,
puts $F$ in its Jacobian ideal. The critical locus is therefore
exactly $D$ as a reduced set. A repeated factor of $F$ would put
a two-dimensional component in this critical locus, which proves
that $F$ is reduced.

The Hessian along the normalization has the exact factorization
\[
 \operatorname{Hess}(F)|_D
 =2t^{21}
 \begin{pmatrix}7t^2\\0\\-5\end{pmatrix}
 \begin{pmatrix}7t^2&0&-5\end{pmatrix},
 \qquad \frac{1}{6}F_{yyy}|_D=-20t^{17}.                       \tag{A.NL.4}
\]
For $t\ne0$, the plane $x=t^5$ is transverse to $D$.
The restricted Hessian has rank one and its radical is the
$y$ direction; its cubic coefficient in that direction is nonzero.
The holomorphic splitting lemma and a local cube root of a unit
give the transverse cusp equation. At the generic point of $D$,
its conormal is locally free and injects into ambient differentials.
Since $F$ and $dF$ vanish there, $F\in I^2\C\{x,y,z\}_I$,
which is precisely $F\in I^{(2)}$.

To decide ordinary-square membership, use the positive weights
$(5,6,7)$. The three generators have weights $12,20,21$.
Among products of two generators, only $g_1^2$ can contribute
to weight $35$: it needs a coefficient of weight $11$, whose
space is spanned by $xy$. Products $g_1g_2$ and $g_1g_3$ would
need weights $3$ and $2$, which do not occur; the other products
already have greater weight. Thus the weight-$35$ part of $I^2$
is exactly $\C\,xyg_1^2$. Its $x^7$ coefficient is zero, whereas
that of $F$ is four. Taking the weight-$35$ component of a
convergent series identity is legitimate because all weights are
positive. This proves $F\notin I^2$ in the analytic local ring.

For comparison of the curve invariants, the normalization conductor
is $t^{10}\C\{t\}$: the semigroup gaps are $1,2,3,4,8,9$.
The $B$-module generated by the differential image, where $B=\C\{t\}$, is $t^4B\,dt$.
The smallest valuation of a two-by-two minor of the displayed
generators' Jacobian matrix is $19$. Thus the normalized conormal
determinant correction is $19-4=15$, while the conductor exponent
is $10$; the complete-intersection adjunction comparison cannot
identify these two quantities here.

The symbolic-square correction is visible directly on the free
normalized conormal. Set $a=dg_1|_D$ and
$b=dg_2|_D+2t^8a$. Then
\[
 \begin{aligned}
 a&=(-t^7,2t^6,-t^5),& b&=(-6t^{15},5t^{14},0),\\
 dg_3|_D&=-3t^9a+tb,&
 q&=t^{-5}(5t^8a-2b)^2.
 \end{aligned}                                                   \tag{A.NL.5}
\]
The first two rows generate the entire normalized conormal image,
and $q$ is the quadratic term, namely half the Hessian.
If $A,B$ are the dual basis, its radical frame $2A+5t^8B$
is represented in the ambient normal quotient by $t^{-6}\partial_y$.
Its intrinsic cubic is therefore $-20t^{-1}$. These are actual
quadratic and cubic poles of orders five and one. The corresponding
combined invariant $q^3k^2$ has pole order $17$; no bound for such
poles on all non-lci curves is deduced from this one computation.
All substitutions, polynomial identities, weighted coefficient
spaces, conormal comparisons, and Jacobian minors are checked
exactly in \path{code/A_non_lci_cusp.py}.
\end{proof}
